% Part: sets-functions-relations
% Chapter: functions
% Section: basics

\documentclass[../../../include/open-logic-section]{subfiles}

\begin{document}

\olfileid{sfr}{fun}{bas}
\olsection{પાયાની બાબતો}

\begin{explain}
\emph{વિધેય} એવું નિરૂપણ છે જે આપેલા ગણના દરેક !!{element}ને
કોઈ (બીજા) આપેલા ગણના ચોક્કસ !!{element} પર મોકલે છે.
દાખલા તરીકે, $1$~ઉમેરવાની ક્રિયા એક વિધેય વ્યાખ્યાયિત કરે છે:
દરેક સંખ્યા~$n$નું નિરૂપણ અનન્ય સંખ્યા~$n+1$ પર થાય છે.

વધુ સામાન્ય રીતે, વિધેયો આગત તરીકે જોડ, ત્રિજોડ વગેરે લઈ
કોઈ પ્રકારનું નિર્ગત આપી શકે છે. પાયાના અંકગણિતમાંથી આપણને
ઘણાં વિધેયો પરિચિત છે. દાખલા તરીકે, સરવાળો અને ગુણાકાર
વિધેયો છે. તેઓ બે સંખ્યાઓ લે છે અને ત્રીજી સંખ્યા આપે છે.

આ ગાણિતિક, અમૂર્ત અર્થમાં વિધેય એક \emph{કાળી પેટી} છે:
કયા આગત સાથે કયું નિર્ગત જોડાયેલું છે એટલું જ મહત્ત્વનું છે,
નિર્ગતની ગણતરી કરવાની પદ્ધતિ નહીં.
\end{explain}

\begin{defn}[વિધેય]
\emph{વિધેય} $f \colon A \to B$ એ $A$ના દરેક !!{element}નું
$B$ના કોઈ !!{element} પરનું નિરૂપણ છે.

આપણે $A$ને $f$નો \emph{પ્રદેશ} અને $B$ને $f$નો
\emph{સહપ્રદેશ} કહીએ છીએ. $A$ના !!{element}sને $f$નાં
આગતો અથવા \emph{દલીલો} કહે છે; અને $f$ દ્વારા દલીલ~$x$ સાથે
જોડાતા $B$ના !!{element}ને દલીલ~$x$ માટેની
\emph{$f$ની કિંમત} કહે છે, જેને $f(x)$ લખાય છે.

$f$નો \emph{વિસ્તાર} $\ran{f}$ એ સહપ્રદેશનો એવો ઉપગણ છે
જેમાં કોઈક દલીલ માટેની $f$ની કિંમતો આવે છે;
$\ran{f} =
\Setabs{f(x)}{x \in A}$.
\end{defn}

\olref{fig:function}માંનું ચિત્ર વિધેયો વિશે વિચારવામાં મદદરૂપ
થઈ શકે છે. ડાબી બાજુનો લંબવૃત્ત વિધેયનો \emph{પ્રદેશ}
દર્શાવે છે; જમણી બાજુનો લંબવૃત્ત તેનો \emph{સહપ્રદેશ}
દર્શાવે છે; અને તીર પ્રદેશની \emph{દલીલ}થી સહપ્રદેશની
તેને અનુરૂપ \emph{કિંમત} તરફ જાય છે.

\begin{figure}
  \olasset{assets/diagrams/function.tikz}
  \caption{વિધેય એ એક ગણના દરેક !!{element}નું બીજા ગણના
    !!a{element} પરનું નિરૂપણ છે. તીર પ્રદેશની દલીલથી
    સહપ્રદેશની તેને અનુરૂપ કિંમત તરફ જાય છે.}
  \ollabel{fig:function}
\end{figure}

\begin{ex}
ગુણાકાર પ્રાકૃતિક સંખ્યાઓની જોડને આગત તરીકે લે છે અને
પ્રાકૃતિક સંખ્યાઓને નિર્ગત તરીકે આપે છે. તેથી તે
$\Nat \times \Nat$ (પ્રદેશ)થી $\Nat$ (સહપ્રદેશ)માં જાય છે.
તેનો વિસ્તાર પણ $\Nat$ છે, કારણ કે દરેક $n \in \Nat$
એ $n \times 1$ છે.
\end{ex}

\begin{ex}
ગુણાકાર વિધેય છે, કારણ કે તે દરેક આગતને---પ્રાકૃતિક સંખ્યાઓની
દરેક જોડને---એક જ નિર્ગત સાથે જોડે છે:
$\times \colon \Nat^2 \to
\Nat$. તેનાથી વિપરીત, પ્રાકૃતિક સંખ્યા~$n$ને
$x^2 = n$ થાય એવી વાસ્તવિક સંખ્યા~$x$ પર મોકલતું નિરૂપણ
વિધેયાત્મક નથી, કારણ કે દરેક ધન પૂર્ણાંક~$n$નાં બે વર્ગમૂળ
છે: $\sqrt{n}$ અને~$-\sqrt{n}$. માત્ર અનૃણ (મુખ્ય) વર્ગમૂળ આપીને
આ નિરૂપણને વિધેયાત્મક બનાવી શકીએ:
$\sqrt{\phantom{X}} \colon
\Nat \to \Real$.
\end{ex}
\sourcecorrection{OLFUN-002}{મૂળની \texttt{function-basics.tex}, પંક્તિઓ
64--71માં પ્રદેશમાં શૂન્યનો સમાવેશ હોવા છતાં ``ધન વર્ગમૂળ'' કહેવાયું છે.
અહીં શૂન્યને પણ આવરી લેતું ``અનૃણ (મુખ્ય) વર્ગમૂળ'' લખ્યું છે.}

\begin{ex}
વર્ગના દરેક વિદ્યાર્થીને તેના અંતિમ ગુણાંક સાથે જોડતો સંબંધ
વિધેય છે---એક જ વર્ગમાં કોઈ વિદ્યાર્થીને બે જુદા અંતિમ
ગુણાંક મળી શકે નહીં. વર્ગના દરેક વિદ્યાર્થીને તેનાં
માતાપિતા સાથે જોડતો સંબંધ વિધેય નથી: વિદ્યાર્થીઓને
શૂન્ય, બે કે વધુ માતાપિતા હોઈ શકે છે.
\end{ex}

\begin{explain}
દરેક સંભવિત દલીલ માટે વિધેયની કિંમત શું છે તે કોઈ ચોક્કસ
રીતે જણાવીને આપણે વિધેયો વ્યાખ્યાયિત કરી શકીએ છીએ.
તે માટે સૂત્ર આપી શકાય, કિંમતની ગણતરી કરવાની પદ્ધતિ
વર્ણવી શકાય, અથવા દરેક દલીલની કિંમતોની યાદી આપી શકાય.
વિધેયને ગમે તે રીતે વ્યાખ્યાયિત કરીએ, દરેક દલીલ માટે
એક અને માત્ર એક કિંમત નક્કી કરી છે તેની ખાતરી કરવી જોઈએ.
\end{explain}

\begin{ex}
$f(x) = x+1$ થાય તે રીતે $f \colon \Nat \to \Nat$
વ્યાખ્યાયિત કરીએ. આ વ્યાખ્યા $f$ને એવું વિધેય બનાવે છે
જે પ્રાકૃતિક સંખ્યાઓ લે છે અને પ્રાકૃતિક સંખ્યાઓ આપે છે.
તે જણાવે છે કે પ્રાકૃતિક સંખ્યા~$x$ આપતાં $f$ તેનું
અનુગામી~$x+1$ આપશે.
આ કિસ્સામાં સહપ્રદેશ $\Nat$ એ $f$નો વિસ્તાર નથી,
કારણ કે પ્રાકૃતિક સંખ્યા~$0$ કોઈ પ્રાકૃતિક સંખ્યાનું
અનુગામી નથી. $f$નો વિસ્તાર બધા ધન પૂર્ણાંકોનો ગણ
$\Int^{+}$ છે.
\end{ex}

\begin{ex}\ollabel{examplefunext}
$g(x) = x+2-1$ થાય તે રીતે $g \colon \Nat \to \Nat$
વ્યાખ્યાયિત કરીએ. આ જણાવે છે કે $g$ પ્રાકૃતિક સંખ્યાઓ
લેતું અને પ્રાકૃતિક સંખ્યાઓ આપતું વિધેય છે.
પ્રાકૃતિક સંખ્યા~$x$ આપતાં $g$, $x$ના અનુગામીના
અનુગામીનો પુરોગામી, એટલે કે $x+1$, આપશે.
\end{ex}
\sourcecorrection{OLFUN-003}{મૂળની \texttt{function-basics.tex}, પંક્તિઓ
103--107માં આગતનું નામ પહેલાં \(n\) અને પછી \(x\) છે. વ્યાખ્યા અને
પછીના સૂત્ર સાથે સુસંગત રહે તે માટે અહીં સર્વત્ર \(x\) વાપર્યું છે.}

\begin{explain}
આપણે હમણાં જુદી \emph{વ્યાખ્યાઓ} ધરાવતા બે વિધેયો,
$f$ અને $g$, જોયાં. જોકે, આ \emph{એક જ વિધેય} છે.
કારણ કે કોઈ પણ પ્રાકૃતિક સંખ્યા~$n$ માટે
$f(n) = n+1 = n+2-1 =
g(n)$ થાય છે. બીજા શબ્દોમાં, $f$ અને~$g$ માટેની આપણી
વ્યાખ્યાઓ જુદાં સમીકરણો દ્વારા એક જ નિરૂપણ નક્કી કરે છે.
આમ, ગર્ભિત રીતે આપણે વિધેયો માટે દલીલોની કિંમતો દ્વારા
નિર્ધારિત સમાનતાના સિદ્ધાંત પર આધાર રાખીએ છીએ:
\[
  \text{જો }\forall x\, f(x) = g(x)\text{, તો }f = g
\]
શરત એટલી કે $f$ અને~$g$નો પ્રદેશ અને સહપ્રદેશ એક જ હોય.
\end{explain}

\begin{ex}
આપણે કિસ્સા પાડી વિધેયો વ્યાખ્યાયિત પણ કરી શકીએ.
દાખલા તરીકે, $h \colon \Nat \to \Nat$ને આમ વ્યાખ્યાયિત
કરી શકીએ:
\[
h(x) =
\begin{cases}
  \frac{x}{2} & \text{જો $x$ યુગ્મ હોય} \\
  \frac{x+1}{2} & \text{જો $x$ અયુગ્મ હોય.}
\end{cases}
\]
દરેક પ્રાકૃતિક સંખ્યા યુગ્મ કે અયુગ્મ હોવાથી, આ વિધેયનું
નિર્ગત હંમેશાં પ્રાકૃતિક સંખ્યા હશે. એટલું યાદ રાખો કે
કિસ્સા પાડી વિધેય વ્યાખ્યાયિત કરો ત્યારે દરેક સંભવિત
આગત ચોક્કસ એક કિસ્સામાં આવવું જોઈએ. કેટલીક વાર
કિસ્સાઓ બધા આગતોને આવરી લે છે અને પરસ્પર અલગ છે
તે સાબિત કરવું જરૂરી બનશે.
\end{ex}

\end{document}
